\begin{figure}[H]
\centering

\begin{tikzpicture}[
>=Latex,
font=\small,
node distance=8mm,
monomer/.style={
draw,
rounded corners,
fill=cyan!20,
minimum width=2.2cm,
minimum height=0.9cm,
align=center
},
polymer/.style={
draw,
rounded corners,
fill=orange!20,
minimum width=10cm,
minimum height=1cm,
align=center
}
]

%%%%%%%%%%%%%%%%%%%%%%%%%%%%%%%%%%%%%%%%%%%%%%%%%%%%%%%%%%%%
% Monómeros
%%%%%%%%%%%%%%%%%%%%%%%%%%%%%%%%%%%%%%%%%%%%%%%%%%%%%%%%%%%%

\node[monomer] (m1) {Monómero};

\node[right=of m1] (plus1) {$+$};

\node[monomer,right=of plus1] (m2) {Monómero};

\node[right=of m2] (plus2) {$+$};

\node[monomer,right=of plus2] (m3) {Monómero};

\node[right=of m3] (dots) {$\cdots$};

%%%%%%%%%%%%%%%%%%%%%%%%%%%%%%%%%%%%%%%%%%%%%%%%%%%%%%%%%%%%
% Flecha
%%%%%%%%%%%%%%%%%%%%%%%%%%%%%%%%%%%%%%%%%%%%%%%%%%%%%%%%%%%%

\draw[->,very thick]
($(m2.south)+(0,-0.2)$)
--
node[right]{\textbf{Polimerización}}
++(0,-1.6);

%%%%%%%%%%%%%%%%%%%%%%%%%%%%%%%%%%%%%%%%%%%%%%%%%%%%%%%%%%%%
% Polímero
%%%%%%%%%%%%%%%%%%%%%%%%%%%%%%%%%%%%%%%%%%%%%%%%%%%%%%%%%%%%

\node[
polymer,
below=2.3cm of m2
] (poly)
{
Monómero -- Monómero -- Monómero -- Monómero -- Monómero -- Monómero
};

%%%%%%%%%%%%%%%%%%%%%%%%%%%%%%%%%%%%%%%%%%%%%%%%%%%%%%%%%%%%
% Texto inferior
%%%%%%%%%%%%%%%%%%%%%%%%%%%%%%%%%%%%%%%%%%%%%%%%%%%%%%%%%%%%

\node[
below=5mm of poly
]
{\textbf{Macromolécula (Polímero)}};

%%%%%%%%%%%%%%%%%%%%%%%%%%%%%%%%%%%%%%%%%%%%%%%%%%%%%%%%%%%%
% Flechas laterales
%%%%%%%%%%%%%%%%%%%%%%%%%%%%%%%%%%%%%%%%%%%%%%%%%%%%%%%%%%%%

\draw[->]
(m1.south)
to[out=-90,in=160]
(poly.west);

\draw[->]
(m3.south)
to[out=-90,in=20]
(poly.east);

%%%%%%%%%%%%%%%%%%%%%%%%%%%%%%%%%%%%%%%%%%%%%%%%%%%%%%%%%%%%

\node[
left=2.2cm of poly,
align=center
]
{Muchos\\monómeros};

\node[
right=2.2cm of poly,
align=center
]
{Cadena\\polimérica};

\end{tikzpicture}

\caption{Formación de un polímero por unión repetitiva de monómeros.}

\end{figure}